\section{问题四：固定轨迹上的时空任务调度}

问题四不再优化机器人路径，而是在问题三的固定融合轨迹上选择“何时、对哪个目标、执行什么任务”。由平滑轨迹计算位置、速度、加速度以及机器人到目标的距离和方位角，并在 0.1 s 候选网格上检查完整准备窗口。

射击候选要求准备窗口内持续满足
\[
5\le d\le30,\qquad v\le2,\qquad a\le1.5,
\]
且准备时间为 1.5 s；拍照候选要求
\[
10\le d\le40,\qquad v\le1.5,\qquad a\le1.5,
\]
并按 0.5 s 对准窗口处理。同一拍照目标的任意两次有效照片满足圆周最小方位角差不小于 $60^\circ$。候选生成后再用 0.01 s 连续状态评估进行复核，避免 10 Hz 离散网格遗漏窗口内部的约束越界。

最终调度采用三阶段字典序 MILP：第一阶段在结果模板九行容量内最大化任务数；第二阶段固定任务数后最大化所有已选任务的最小归一化硬约束裕度；第三阶段保持前两级最优值，进一步最大化总裕度与拍照角度多样性。正式方案完成 \QFourTaskCount 项任务，其中射击 \QFourShootCount 项、拍照 \QFourPhotoCount 项，射击期望命中数为 \QFourExpectedHits。MILP 最小裕度为 \QFourMinMargin，而最早可行贪心基线仅为 \QFourGreedyMargin，前者约为后者的 \QFourMarginRatio 倍。

题面未明确说明不同任务的准备过程是否允许并行。正式主模型采用“单机器人同一时段只准备/执行一项任务”的保守工程假设，但该假设不作为题面原文引用。仓库同时运行允许准备阶段并行的对照模型，并将两种资源模式的一级最优任务数与安全裕度写入独立结果文件，从而判断该附加假设是否改变核心结论。

\input{generated/table_q4_summary.tex}
\input{generated/table_q4_schedule.tex}
